\subsection{Filling Out the Design Space: Cartesian Base}
\label{sec:cartesian-base}

A central idea in this work has been to consider the `base' notion of usage in
graded systems: linear-base systems have linear propositions, with a linear
axiom rule, and linear function type; graded-base systems have graded propositions,
with a graded axiom rule, and a graded function type. This framing suggest
an unexplored point in the design space: what about a graded system with
a cartesian base, i.e., with a cartesian axiom rule and standard cartesian
function (type), with contexts comprising both cartesian and graded
assumptions?

Crucially we must then consider how the cartesian part interacts with
the graded part. On the left, we have the rules for dereliction and
application in Linear Base, on the right have the potential rules
for a Cartesian Base. Since cartesian propositions can be used arbitrarily
we now require an element $\omega \in \mathcal{R}$ of the semiring which
behaves in an absorbing fashion for both addition $+$ and multiplication $\cdot$.
Dereliction in Cartesian Base can then mark a cartesian assumption
with $\omega$. Application must also scale the argument context $\Gamma_2$
by $\omega$ since the function arrow is cartesian and so may use
its argument arbitrarily:

\begin{minipage}{0.5\linewidth}
\begin{align*}
\begin{array}{c}
\text{(Linear Base)} \\[0.5em]
\dfrac{\Gamma, x : A \vdash_{\textsc{l}} t : B}
      {\Gamma, x : [A]_1 \vdash_{\textsc{l}} t : B} (der)
\\[1.5em]
\dfrac{\Gamma_1 \vdash_{\textsc{l}} t_1 : A \multimap B \quad \Gamma_2 \vdash_{\textsc{l}} t_2 : A}
      {\Gamma_1 + \Gamma_2 \vdash_{\textsc{l}} t_1\ t_2 : B}(app)
\end{array}
\\
\end{align*}
\end{minipage}
\begin{minipage}{0.5\linewidth}
\begin{align*}
\begin{array}{c}
\text{(Cartesian Base)} \\[0.5em]
\dfrac{\Gamma, x : A \vdash t : B}
      {\Gamma, x : [A]_\omega \vdash t : B}(der)
\\[1.5em]
\dfrac{\Gamma_1 \vdash t_1 : A \rightarrow B \quad \Gamma_2 \vdash t_2 : A}
      {\Gamma_1 + \omega \ast \Gamma_2 \vdash t_1\ t_2 : B}(app)
\end{array}
\\
\end{align*}
\end{minipage}

\noindent
The above seems a reasonable approach to building a
Cartesian Base system. However, this system
provides little benefit as the grading quickly degenerates to
being always $\omega$ in derivations. Consider below the derivation of
the `copy' operation on an axiom to produce a pair:

\begin{align*}
\begin{array}{cc}
\dfrac{
\dfrac{x : A \vdash_{\textsc{l}} x : A}
      {x : [A]_1 \vdash_{\textsc{l}} x : A}
\quad
\dfrac{x : A \vdash_{\textsc{l}} x : A}
      {x : [A]_1 \vdash_{\textsc{l}} x : A}
}
{x : [A]_{1 + 1} \vdash_{\textsc{l}} (x, x) : A \otimes A}
&
\dfrac{
\dfrac{x : A \vdash x : A}
      {x : [A]_\omega \vdash x : A}
\quad
\dfrac{x : A \vdash x : A}
      {x : [A]_\omega \vdash x : A}
}
{x : [A]_\omega \vdash (x, x) : A \otimes A}
\end{array}
\end{align*}
%
In Linear Base, there is a precise account of the non-linearity
incurred: the term uses $x$ with graded $1+1$. However, in
Cartesian Base we can only say $\omega$ due to the dereliction rule.
In the end, the Cartesian Base system
is ultimately no more informative as a graded system than
standard intuitionistic, cartesian systems.
